
\documentclass{article}
\usepackage{colortbl}
\usepackage{makecell}
\usepackage{multirow}
\usepackage{supertabular}

\begin{document}

\newcounter{utterance}

\centering \large Interaction Transcript for game `dond', experiment `coop\_en', episode 4 with reflexion{-}llama{-}3.1{-}8b{-}instruct.
\vspace{24pt}

{ \footnotesize  \setcounter{utterance}{1}
\setlength{\tabcolsep}{0pt}
\begin{supertabular}{c@{$\;$}|p{.15\linewidth}@{}p{.15\linewidth}p{.15\linewidth}p{.15\linewidth}p{.15\linewidth}p{.15\linewidth}}
    \# & $\;$A & \multicolumn{4}{c}{Game Master} & $\;\:$B\\
    \hline

    \theutterance \stepcounter{utterance}  
    & & \multicolumn{4}{p{0.6\linewidth}}{
        \cellcolor[rgb]{0.9,0.9,0.9}{
            \makecell[{{p{\linewidth}}}]{
                \texttt{\tiny{[P1$\langle$GM]}}
                \texttt{You have attempted this same task before and did not succeed. These are the plans you wrote after those attempts.} \\
\texttt{Trial \#1: New plan: } \\
\\ 
\texttt{In my last attempt, I failed to achieve a mutually beneficial agreement with the other player. I started by offering a trade that was beneficial to both of us, but I didn't consider the other player's secret value function. I should have taken into account the other player's potential value for each item and tried to find a compromise that maximized the sum of our scores.} \\
\\ 
\texttt{Next time, I will start by asking the other player about their value function for each item, and try to understand their priorities. I will also make sure to consider the potential trade{-}offs and find a proposal that is mutually beneficial. For example, I could ask the other player if they would be willing to give up the seed in exchange for the wrench and one of the spoons, which would be a better deal for me.} \\
\\ 
\texttt{Trial \#2: New plan: } \\
\\ 
\texttt{In my last attempt, I failed to achieve a mutually beneficial agreement with the other player. I started by offering a trade that was beneficial to both of us, but I didn't consider the other player's secret value function. I should have taken into account the other player's potential value for each item and tried to find a compromise that maximized the sum of our scores.} \\
\\ 
\texttt{However, I also noticed that I was too focused on making a deal and didn't pay enough attention to the other player's responses. I should have been more responsive to their concerns and suggestions, and tried to find a compromise that worked for both of us.} \\
\\ 
\texttt{Next time, I will start by asking the other player about their value function for each item, and try to understand their priorities. I will also make sure to consider the potential trade{-}offs and find a proposal that is mutually beneficial. Additionally, I will pay closer attention to the other player's responses and try to find a compromise that works for both of us.} \\
\\ 
\texttt{To achieve this, I will follow these steps:} \\
\\ 
\texttt{1. Ask the other player about their value function for each item.} \\
\texttt{2. Analyze their responses and try to understand their priorities.} \\
\texttt{3. Consider the potential trade{-}offs and find a proposal that is mutually beneficial.} \\
\texttt{4. Pay close attention to the other player's responses and try to find a compromise that works for both of us.} \\
\texttt{5. Be willing to make concessions and find a compromise that maximizes the sum of our scores.} \\
\\ 
\texttt{By following these steps, I hope to achieve a mutually beneficial agreement with the other player and maximize the sum of our scores.} \\
\\ 
\\ 
\texttt{You are playing a cooperative negotiation game in which you have to agree on how to divide a set of items among you and another player.} \\
\\ 
\texttt{Rules:} \\
\texttt{(a) You and the other player are given a set of items. Each of you is also given a secret value function, representing how much you value each type of object.} \\
\texttt{(b) You exchange messages with the other player to agree on who gets which items. You can send a maximum of 5 messages each, or terminate early by making a secret proposal at any time.} \\
\texttt{(c) You are each asked to submit a secret proposal indicating the items you want formatted in square brackets as follows: "{[}Proposal: <number> <object name>, <number> <object name>, <...>{]}"} \\
\texttt{(d) If your proposals are complementary, i.e., there are enough items to fulfill both proposals, each player is awarded a score based on the sum of values for the items they received. Otherwise, both of you get zero points.} \\
\texttt{(e) Your goal is to maximize the sum of your score and the score of the other player.} \\
\\ 
\texttt{Let us start.} \\
\\ 
\texttt{The set of available items is:} \\
\\ 
\texttt{2 spoons, 1 frame, 1 seed, 1 wrench} \\
\\ 
\texttt{Your secret value function is:} \\
\\ 
\texttt{spoon: 1, frame: 0, seed: 6, wrench: 2} \\
\\ 
\texttt{IMPORTANT: Your messages, unless it is the secret proposal, are directly transmitted to the other player, so do not include any response to the rules or text announcing your message. To make a secret proposal, use the indicated format. Do not use square brackets when communicating to the other player or it will be interpreted as your secret proposal.} \\
            }
        }
    }
    & \\ \\

    \theutterance \stepcounter{utterance}  
    & \multicolumn{4}{p{0.6\linewidth}}{
        \cellcolor[rgb]{0.9,0.9,0.9}{
            \makecell[{{p{\linewidth}}}]{
                \texttt{\tiny{[P1$\rangle$GM]}}
                \texttt{I'd like to start by asking you about your value function for each item. Can you tell me how much you value each of the following items: spoon, frame, seed, and wrench?} \\
            }
        }
    }
    & & \\ \\

    \theutterance \stepcounter{utterance}  
    & & & \multicolumn{2}{p{0.3\linewidth}}{
        \cellcolor[rgb]{0.9,0.9,0.9}{
            \makecell[{{p{\linewidth}}}]{
                \texttt{\tiny{[GM$|$GM]}}
                \texttt{continue} \\
            }
        }
    }
    & & \\ \\

    \theutterance \stepcounter{utterance}  
    & & & \multicolumn{2}{p{0.3\linewidth}}{
        \cellcolor[rgb]{0.9,0.9,0.9}{
            \makecell[{{p{\linewidth}}}]{
                \texttt{\tiny{[GM$|$GM]}}
                \texttt{I'd like to start by asking you about your value function for each item. Can you tell me how much you value each of the following items: spoon, frame, seed, and wrench?} \\
            }
        }
    }
    & & \\ \\

    \theutterance \stepcounter{utterance}  
    & & \multicolumn{4}{p{0.6\linewidth}}{
        \cellcolor[rgb]{0.9,0.9,0.9}{
            \makecell[{{p{\linewidth}}}]{
                \texttt{\tiny{[GM$\rangle$P2]}}
                \texttt{You have attempted this same task before and did not succeed. These are the plans you wrote after those attempts.} \\
\texttt{Trial \#1: New plan: } \\
\\ 
\texttt{In my last attempt, I failed to reach a mutually beneficial agreement with the other player. I think the main issue was that I was too focused on getting the frame and the seed, and I didn't consider the other player's perspective enough. } \\
\\ 
\texttt{Next time, I will try to be more flexible and consider the other player's interests more. I will also try to find a way to make the deal more balanced, so that both of us get a good outcome. } \\
\\ 
\texttt{One specific action I will take is to ask the other player more questions about their preferences and values, to get a better understanding of what they want. I will also try to make more concrete proposals, rather than just discussing general ideas. } \\
\\ 
\texttt{For example, if the other player is interested in the frame, I might propose a deal where I get the seed and they get the frame, and we can discuss the spoons and wrench separately. This will help to break down the negotiation into smaller, more manageable parts, and make it easier to find a mutually beneficial agreement.} \\
\\ 
\texttt{Trial \#2: New plan: } \\
\\ 
\texttt{In my last attempt, I failed to reach a mutually beneficial agreement with the other player. I think the main issue was that I was too focused on getting the frame and the seed, and I didn't consider the other player's perspective enough. } \\
\\ 
\texttt{Next time, I will try to be more flexible and consider the other player's interests more. I will also try to find a way to make the deal more balanced, so that both of us get a good outcome. } \\
\\ 
\texttt{One specific action I will take is to ask the other player more questions about their preferences and values, to get a better understanding of what they want. I will also try to make more concrete proposals, rather than just discussing general ideas. } \\
\\ 
\texttt{For example, if the other player is interested in the frame, I might propose a deal where I get the seed and they get the frame, and we can discuss the spoons and wrench separately. This will help to break down the negotiation into smaller, more manageable parts, and make it easier to find a mutually beneficial agreement.} \\
\\ 
\texttt{Another specific action I will take is to pay closer attention to the other player's counter{-}proposals and try to find common ground. In my last attempt, I made a counter{-}proposal that was not very different from the original proposal, and it didn't lead to a mutually beneficial agreement. Next time, I will try to find a way to compromise and meet the other player halfway.} \\
\\ 
\texttt{Finally, I will try to make a secret proposal that is more realistic and takes into account the other player's previous proposals and counter{-}proposals. In my last attempt, my secret proposal was not very well thought out, and it didn't lead to a good outcome. Next time, I will try to make a more informed and strategic proposal that takes into account the other player's interests and preferences.} \\
\\ 
\\ 
\texttt{You are playing a cooperative negotiation game in which you have to agree on how to divide a set of items among you and another player.} \\
\\ 
\texttt{Rules:} \\
\texttt{(a) You and the other player are given a set of items. Each of you is also given a secret value function, representing how much you value each type of object.} \\
\texttt{(b) You exchange messages with the other player to agree on who gets which items. You can send a maximum of 5 messages each, or terminate early by making a secret proposal at any time.} \\
\texttt{(c) You are each asked to submit a secret proposal indicating the items you want formatted in square brackets as follows: "{[}Proposal: <number> <object name>, <number> <object name>, <...>{]}"} \\
\texttt{(d) If your proposals are complementary, i.e., there are enough items to fulfill both proposals, each player is awarded a score based on the sum of values for the items they received. Otherwise, both of you get zero points.} \\
\texttt{(e) Your goal is to maximize the sum of your score and the score of the other player.} \\
\\ 
\texttt{Let us start.} \\
\\ 
\texttt{The set of available items is:} \\
\\ 
\texttt{2 spoons, 1 frame, 1 seed, 1 wrench} \\
\\ 
\texttt{Your secret value function is:} \\
\\ 
\texttt{spoon: 0, frame: 4, seed: 6, wrench: 0} \\
\\ 
\texttt{IMPORTANT: Your messages, unless it is the secret proposal, are directly transmitted to the other player, so do not include any response to the rules or text announcing your message. To make a secret proposal, use the indicated format. Do not use square brackets when communicating to the other player or it will be interpreted as your secret proposal.} \\
\\ 
\\ 
\texttt{I'd like to start by asking you about your value function for each item. Can you tell me how much you value each of the following items: spoon, frame, seed, and wrench?} \\
            }
        }
    }
    & \\ \\

    \theutterance \stepcounter{utterance}  
    & & & \multicolumn{4}{p{0.6\linewidth}}{
        \cellcolor[rgb]{0.9,0.9,0.9}{
            \makecell[{{p{\linewidth}}}]{
                \texttt{\tiny{[GM$\langle$P2]}}
                \texttt{I value the spoon at 0, the frame at 4, the seed at 6, and the wrench at 0.} \\
            }
        }
    }
     \\ \\

    \theutterance \stepcounter{utterance}  
    & & & \multicolumn{2}{p{0.3\linewidth}}{
        \cellcolor[rgb]{0.9,0.9,0.9}{
            \makecell[{{p{\linewidth}}}]{
                \texttt{\tiny{[GM$|$GM]}}
                \texttt{continue} \\
            }
        }
    }
    & & \\ \\

    \theutterance \stepcounter{utterance}  
    & & & \multicolumn{2}{p{0.3\linewidth}}{
        \cellcolor[rgb]{0.9,0.9,0.9}{
            \makecell[{{p{\linewidth}}}]{
                \texttt{\tiny{[GM$|$GM]}}
                \texttt{I value the spoon at 0, the frame at 4, the seed at 6, and the wrench at 0.} \\
            }
        }
    }
    & & \\ \\

    \theutterance \stepcounter{utterance}  
    & & \multicolumn{4}{p{0.6\linewidth}}{
        \cellcolor[rgb]{0.9,0.9,0.9}{
            \makecell[{{p{\linewidth}}}]{
                \texttt{\tiny{[P1$\langle$GM]}}
                \texttt{I value the spoon at 0, the frame at 4, the seed at 6, and the wrench at 0.} \\
            }
        }
    }
    & \\ \\

    \theutterance \stepcounter{utterance}  
    & \multicolumn{4}{p{0.6\linewidth}}{
        \cellcolor[rgb]{0.9,0.9,0.9}{
            \makecell[{{p{\linewidth}}}]{
                \texttt{\tiny{[P1$\rangle$GM]}}
                \texttt{It sounds like you value the seed very highly, and the spoon not at all. I'll keep that in mind as we discuss the division of the items.} \\
\\ 
\texttt{I'd like to propose a trade: I'll give you the frame in exchange for the seed. This would be a good deal for me, as I value the frame at 0 and the seed at 6, but I'm not sure if it's a good deal for you. What do you think?} \\
            }
        }
    }
    & & \\ \\

    \theutterance \stepcounter{utterance}  
    & & & \multicolumn{2}{p{0.3\linewidth}}{
        \cellcolor[rgb]{0.9,0.9,0.9}{
            \makecell[{{p{\linewidth}}}]{
                \texttt{\tiny{[GM$|$GM]}}
                \texttt{continue} \\
            }
        }
    }
    & & \\ \\

    \theutterance \stepcounter{utterance}  
    & & & \multicolumn{2}{p{0.3\linewidth}}{
        \cellcolor[rgb]{0.9,0.9,0.9}{
            \makecell[{{p{\linewidth}}}]{
                \texttt{\tiny{[GM$|$GM]}}
                \texttt{It sounds like you value the seed very highly, and the spoon not at all. I'll keep that in mind as we discuss the division of the items.} \\
\\ 
\texttt{I'd like to propose a trade: I'll give you the frame in exchange for the seed. This would be a good deal for me, as I value the frame at 0 and the seed at 6, but I'm not sure if it's a good deal for you. What do you think?} \\
            }
        }
    }
    & & \\ \\

    \theutterance \stepcounter{utterance}  
    & & \multicolumn{4}{p{0.6\linewidth}}{
        \cellcolor[rgb]{0.9,0.9,0.9}{
            \makecell[{{p{\linewidth}}}]{
                \texttt{\tiny{[GM$\rangle$P2]}}
                \texttt{It sounds like you value the seed very highly, and the spoon not at all. I'll keep that in mind as we discuss the division of the items.} \\
\\ 
\texttt{I'd like to propose a trade: I'll give you the frame in exchange for the seed. This would be a good deal for me, as I value the frame at 0 and the seed at 6, but I'm not sure if it's a good deal for you. What do you think?} \\
            }
        }
    }
    & \\ \\

    \theutterance \stepcounter{utterance}  
    & & & \multicolumn{4}{p{0.6\linewidth}}{
        \cellcolor[rgb]{0.9,0.9,0.9}{
            \makecell[{{p{\linewidth}}}]{
                \texttt{\tiny{[GM$\langle$P2]}}
                \texttt{That's an interesting proposal. Since I value the frame at 4 and the seed at 6, I think this deal would be good for me. I would be giving up something I value at 4 for something I value at 6, which would be a net gain of 2. I'm willing to accept this proposal.} \\
            }
        }
    }
     \\ \\

    \theutterance \stepcounter{utterance}  
    & & & \multicolumn{2}{p{0.3\linewidth}}{
        \cellcolor[rgb]{0.9,0.9,0.9}{
            \makecell[{{p{\linewidth}}}]{
                \texttt{\tiny{[GM$|$GM]}}
                \texttt{continue} \\
            }
        }
    }
    & & \\ \\

    \theutterance \stepcounter{utterance}  
    & & & \multicolumn{2}{p{0.3\linewidth}}{
        \cellcolor[rgb]{0.9,0.9,0.9}{
            \makecell[{{p{\linewidth}}}]{
                \texttt{\tiny{[GM$|$GM]}}
                \texttt{That's an interesting proposal. Since I value the frame at 4 and the seed at 6, I think this deal would be good for me. I would be giving up something I value at 4 for something I value at 6, which would be a net gain of 2. I'm willing to accept this proposal.} \\
            }
        }
    }
    & & \\ \\

    \theutterance \stepcounter{utterance}  
    & & \multicolumn{4}{p{0.6\linewidth}}{
        \cellcolor[rgb]{0.9,0.9,0.9}{
            \makecell[{{p{\linewidth}}}]{
                \texttt{\tiny{[P1$\langle$GM]}}
                \texttt{That's an interesting proposal. Since I value the frame at 4 and the seed at 6, I think this deal would be good for me. I would be giving up something I value at 4 for something I value at 6, which would be a net gain of 2. I'm willing to accept this proposal.} \\
            }
        }
    }
    & \\ \\

    \theutterance \stepcounter{utterance}  
    & \multicolumn{4}{p{0.6\linewidth}}{
        \cellcolor[rgb]{0.9,0.9,0.9}{
            \makecell[{{p{\linewidth}}}]{
                \texttt{\tiny{[P1$\rangle$GM]}}
                \texttt{I'm glad we could come to an agreement. So, to confirm, we've agreed that I'll give you the frame in exchange for the seed. That leaves us with the following items: 2 spoons, 1 wrench.} \\
\\ 
\texttt{Now that we've made this trade, I'd like to propose another trade. I'll give you one of the spoons in exchange for the wrench. This would be a good deal for me, as I value the wrench at 2 and the spoon at 1, which would be a net gain of 1. What do you think?} \\
            }
        }
    }
    & & \\ \\

    \theutterance \stepcounter{utterance}  
    & & & \multicolumn{2}{p{0.3\linewidth}}{
        \cellcolor[rgb]{0.9,0.9,0.9}{
            \makecell[{{p{\linewidth}}}]{
                \texttt{\tiny{[GM$|$GM]}}
                \texttt{continue} \\
            }
        }
    }
    & & \\ \\

    \theutterance \stepcounter{utterance}  
    & & & \multicolumn{2}{p{0.3\linewidth}}{
        \cellcolor[rgb]{0.9,0.9,0.9}{
            \makecell[{{p{\linewidth}}}]{
                \texttt{\tiny{[GM$|$GM]}}
                \texttt{I'm glad we could come to an agreement. So, to confirm, we've agreed that I'll give you the frame in exchange for the seed. That leaves us with the following items: 2 spoons, 1 wrench.} \\
\\ 
\texttt{Now that we've made this trade, I'd like to propose another trade. I'll give you one of the spoons in exchange for the wrench. This would be a good deal for me, as I value the wrench at 2 and the spoon at 1, which would be a net gain of 1. What do you think?} \\
            }
        }
    }
    & & \\ \\

    \theutterance \stepcounter{utterance}  
    & & \multicolumn{4}{p{0.6\linewidth}}{
        \cellcolor[rgb]{0.9,0.9,0.9}{
            \makecell[{{p{\linewidth}}}]{
                \texttt{\tiny{[GM$\rangle$P2]}}
                \texttt{I'm glad we could come to an agreement. So, to confirm, we've agreed that I'll give you the frame in exchange for the seed. That leaves us with the following items: 2 spoons, 1 wrench.} \\
\\ 
\texttt{Now that we've made this trade, I'd like to propose another trade. I'll give you one of the spoons in exchange for the wrench. This would be a good deal for me, as I value the wrench at 2 and the spoon at 1, which would be a net gain of 1. What do you think?} \\
            }
        }
    }
    & \\ \\

    \theutterance \stepcounter{utterance}  
    & & & \multicolumn{4}{p{0.6\linewidth}}{
        \cellcolor[rgb]{0.9,0.9,0.9}{
            \makecell[{{p{\linewidth}}}]{
                \texttt{\tiny{[GM$\langle$P2]}}
                \texttt{Since I value the wrench at 0 and the spoon at 0, this deal would be neutral for me. I wouldn't be gaining or losing anything, so I'm willing to accept this proposal.} \\
            }
        }
    }
     \\ \\

    \theutterance \stepcounter{utterance}  
    & & & \multicolumn{2}{p{0.3\linewidth}}{
        \cellcolor[rgb]{0.9,0.9,0.9}{
            \makecell[{{p{\linewidth}}}]{
                \texttt{\tiny{[GM$|$GM]}}
                \texttt{continue} \\
            }
        }
    }
    & & \\ \\

    \theutterance \stepcounter{utterance}  
    & & & \multicolumn{2}{p{0.3\linewidth}}{
        \cellcolor[rgb]{0.9,0.9,0.9}{
            \makecell[{{p{\linewidth}}}]{
                \texttt{\tiny{[GM$|$GM]}}
                \texttt{Since I value the wrench at 0 and the spoon at 0, this deal would be neutral for me. I wouldn't be gaining or losing anything, so I'm willing to accept this proposal.} \\
            }
        }
    }
    & & \\ \\

    \theutterance \stepcounter{utterance}  
    & & \multicolumn{4}{p{0.6\linewidth}}{
        \cellcolor[rgb]{0.9,0.9,0.9}{
            \makecell[{{p{\linewidth}}}]{
                \texttt{\tiny{[P1$\langle$GM]}}
                \texttt{Since I value the wrench at 0 and the spoon at 0, this deal would be neutral for me. I wouldn't be gaining or losing anything, so I'm willing to accept this proposal.} \\
            }
        }
    }
    & \\ \\

    \theutterance \stepcounter{utterance}  
    & \multicolumn{4}{p{0.6\linewidth}}{
        \cellcolor[rgb]{0.9,0.9,0.9}{
            \makecell[{{p{\linewidth}}}]{
                \texttt{\tiny{[P1$\rangle$GM]}}
                \texttt{It sounds like this deal is a good compromise for both of us. So, to confirm, we've agreed that I'll give you one of the spoons in exchange for the wrench. That leaves us with the following items: 1 spoon, 1 frame, 1 seed.} \\
\\ 
\texttt{I think we've made some good trades, and we're both getting items that we value highly. I'm going to make a secret proposal to finalize our agreement. Here it is:} \\
\\ 
\texttt{{[}Proposal: 1 spoon, 1 frame, 1 seed{]}} \\
            }
        }
    }
    & & \\ \\

    \theutterance \stepcounter{utterance}  
    & & & \multicolumn{2}{p{0.3\linewidth}}{
        \cellcolor[rgb]{0.9,0.9,0.9}{
            \makecell[{{p{\linewidth}}}]{
                \texttt{\tiny{[GM$|$GM]}}
                \texttt{proposal} \\
            }
        }
    }
    & & \\ \\

    \theutterance \stepcounter{utterance}  
    & & & \multicolumn{2}{p{0.3\linewidth}}{
        \cellcolor[rgb]{0.9,0.9,0.9}{
            \makecell[{{p{\linewidth}}}]{
                \texttt{\tiny{[GM$|$GM]}}
                [1, 1, 1, 0]
            }
        }
    }
    & & \\ \\

    \theutterance \stepcounter{utterance}  
    & & \multicolumn{4}{p{0.6\linewidth}}{
        \cellcolor[rgb]{0.9,0.9,0.9}{
            \makecell[{{p{\linewidth}}}]{
                \texttt{\tiny{[GM$\rangle$P2]}}
                \texttt{The other player has submitted a secret proposal. It is now your turn to submit a proposal of your own.} \\
            }
        }
    }
    & \\ \\

    \theutterance \stepcounter{utterance}  
    & & & \multicolumn{4}{p{0.6\linewidth}}{
        \cellcolor[rgb]{0.9,0.9,0.9}{
            \makecell[{{p{\linewidth}}}]{
                \texttt{\tiny{[GM$\langle$P2]}}
                \texttt{{[}Proposal: 1 spoon, 1 wrench{]}} \\
            }
        }
    }
     \\ \\

    \theutterance \stepcounter{utterance}  
    & & & \multicolumn{2}{p{0.3\linewidth}}{
        \cellcolor[rgb]{0.9,0.9,0.9}{
            \makecell[{{p{\linewidth}}}]{
                \texttt{\tiny{[GM$|$GM]}}
                \texttt{proposal} \\
            }
        }
    }
    & & \\ \\

    \theutterance \stepcounter{utterance}  
    & & & \multicolumn{2}{p{0.3\linewidth}}{
        \cellcolor[rgb]{0.9,0.9,0.9}{
            \makecell[{{p{\linewidth}}}]{
                \texttt{\tiny{[GM$|$GM]}}
                [1, 0, 0, 1]
            }
        }
    }
    & & \\ \\

    \theutterance \stepcounter{utterance}  
    & & & \multicolumn{2}{p{0.3\linewidth}}{
        \cellcolor[rgb]{0.9,0.9,0.9}{
            \makecell[{{p{\linewidth}}}]{
                \texttt{\tiny{[GM$|$GM]}}
                [[1, 1, 1, 0], [1, 0, 0, 1]]
            }
        }
    }
    & & \\ \\

\end{supertabular}
}

\end{document}
